\emph{(i)} An antilinear map commuting with the face action is $A \circ \overline{\phantom{v}}$ with $A$ commuting with every permutation matrix $P$ (the sign of the twist $-P$ cancels on both sides), since the permutation matrices are real. The face representation is $\mathrm{sgn} \oplus \mathbf{2}$ up to the twist, two inequivalent irreducibles each of multiplicity one, so by Schur's lemma $A = xP_{s} + yP_{d}$ with $x, y \in \mathrm{span}\{1, I_4\}$. As $P_{s}$, $P_{d}$ are real projectors, $\sigma^{2}v = A\bar{A}v = (|x|^{2}P_{s} + |y|^{2}P_{d})v$, so $\sigma$ is an involution exactly when $|x| = |y| = 1$. Compatibility: $C(\sigma v) = \bar{v}^{\mathsf{T}}A^{\mathsf{T}}A\bar{v}$ and $A^{\mathsf{T}}A = x^{2}P_{s} + y^{2}P_{d}$, which equals $1$ exactly when $x^{2} = y^{2} = 1$. The fixed slice of $\sigma_{1,-1}$ is $\{tu + I_4 d : t \in \R,\ d \text{ real doublet}\}$ with $u = u'/\sqrt3$, on which $C = t^{2} - Q(d)$ since $C(u, d) = 0$; that of $\sigma_{-1,1}$ is $\{I_4 tu + d\}$, on which $C = -(t^{2} - Q(d))$; the other two have $C = \pm(t^{2} + Q(d))$, Euclidean.

\emph{(ii)} $P_{s}(e v) = e\kappa u'$ and $P_{d}(e v) = e\kappa w f_{+}$ for nonzero $e \in \mathrm{span}\{1, I_4\}$, since $f_{+} \perp u'$ ($\sum_{k}\zeta^{k} = 0$). If $ev$ were fixed by $\sigma_{x,y}$, its doublet part $d$ would satisfy $d = y\bar{d}$, so $d_{k}\bar{d}_{0} = y\bar{d}_{k}\bar{d}_{0}$ and, with $d_{0} = e\kappa w \neq 0$ (as $\lambda > 0$), $d_{1}/d_{0} = \overline{d_{1}/d_{0}}$ would be real. But $d_{1}/d_{0} = \zeta \notin \R$. The line $\mathrm{span}\{1, I_4\}\,f_{+}$ is $C$-null by Proposition \claimref{families.chart-computation-sign-twist-bracket-form}(2), which also records that it contains no nonzero real vector; the argument above is the stronger statement that it contains no complex multiple of one.

\emph{(iii)} The doublet part of $\tfrac12(v + \sigma_{x,y}v)$ is $\tfrac12(D + y\bar{D})$ with $D = \kappa w f_{+}$. Writing $y = \eta^{2}$, $|\eta| = 1$, this is $\eta\,\mathrm{Re}(\bar\eta D)$, of Euclidean square $\sum_{k}\mathrm{Re}(\bar\eta D_{k})^{2} = \tfrac12\bigl(\|D\|^{2} + \mathrm{Re}(\bar\eta^{2}C(D))\bigr) = \tfrac12\|D\|^{2}$, because $C(D) = \kappa^{2}w^{2}C(f_{+}) = 0$; and $\|D\|^{2} = |\kappa|^{2}\lambda^{2}\|f_{+}\|^{2} = 3|\kappa|^{2}\lambda^{2}$. For $\kappa = -I_4|\kappa|$ the singlet $\kappa u'$ is fixed exactly when $x\bar\kappa = \kappa$, that is $x = -1$; among the $C$-compatible structures that leaves $\sigma_{-1,\mp1}$. For $\sigma_{-1,-1}$ the slice is $I_4\R^{3}$ and the real coordinate is $I_4 p$: $I_4 p_{k} = |\kappa|\bigl(1 + \mathrm{Re}(w\zeta^{k})\bigr) = |\kappa|\bigl(1 + \lambda\cos(\chi + 2\pi k/3)\bigr)$. For $\sigma_{-1,1}$ the doublet is kept real and the singlet imaginary, and the real coordinates are $|\kappa|\bigl(1 + \mathrm{Re}(-I_4 w\zeta^{k})\bigr) = |\kappa|\bigl(1 + \lambda\sin(\chi + 2\pi k/3)\bigr) = |\kappa|\bigl(1 + \lambda\cos(\chi - \pi/2 + 2\pi k/3)\bigr)$, with the singlet coordinate oriented by $-I_4 \mapsto +1$; orienting it by $-I_4 \mapsto -1$ multiplies the real coordinates by $-1$ on the singlet only, which, after an overall sign, is $\chi \mapsto \chi + \pi$ on the doublet, that is $\delta = \chi + \pi/2$; $Q$ is unchanged. The restriction of $C$ is (i).

\emph{(iv)} By Proposition \claimref{families.koide-nullity-condition-one-member-invariant}, $Q_{p} = |p|^{2}/(3s^{2})$ with $s = p \cdot u$. Here $s^{2} = 3|\kappa|^{2}$ and, by (iii), $|p|^{2} = s^{2} + Q(d) = 3|\kappa|^{2} + \tfrac32|\kappa|^{2}\lambda^{2}$, so $Q_{p} = \tfrac13 + \tfrac{\lambda^{2}}{6}$, and $Q_{p} = \tfrac23 \iff \lambda^{2} = 2 \iff s^{2} = Q(d)$, the nullity of that Proposition. The equivalent forms: $C(v) = \kappa^{2}\bigl(C(u') + 2wC(u', f_{+}) + w^{2}C(f_{+})\bigr) = 3\kappa^{2}$ and $\|v\|^{2} = 3|\kappa|^{2}(1 + \lambda^{2})$, so $\|v\|^{2} = 3|C(v)| \iff \lambda^{2} = 2$; the Hermitian weights are $\|P_{s}v\|^{2} = 3|\kappa|^{2}$ and $\|P_{d}v\|^{2} = 3|\kappa|^{2}\lambda^{2}$, in ratio $1:2$ exactly then. For the pairing vector of Proposition \claimref{families.rigidity-pairing-sector-one-event}, $a = a_{A} + \lambda e^{I_4\chi}a_{B}$ with $a_{A} = -I_4 u'/\sqrt3$ and $a_{B} = -I_4 f_{+}/\sqrt3$: $a \cdot a = -\tfrac13\bigl(3 + 0 + 0\bigr) = -1$ and $|a|^{2} = 1 + \lambda^{2}$, so $|a \cdot a| = \tfrac13|a|^{2} \iff \lambda^{2} = 2$. At $\lambda = 1$, $Q_{p} = \tfrac12$.

The script verifies, in exact arithmetic in $\Q(\sqrt3)[I_4]$: the isotropy of the Fourier lines and the half-norm of every real slice of $f_{+}$ at rational points of the unit circle; that each $\sigma_{x,y}$ at those points is an equivariant involution, compatible with $C$ exactly at $x, y = \pm1$, the signature of $C$ on a rational basis of each of the four compatible slices ($(3,0)$, $(1,2)$, $(2,1)$, $(0,3)$ for $(x,y) = (1,1), (1,-1), (-1,1), (-1,-1)$), and fixes no multiple of $v$ for seven values of $w$; the projections, their real coordinates and the value of $Q_{p}$ for those $w$; and the Koide instance $w = 1 + I_4$ ($\lambda^{2} = 2$, $\chi = \pi/4$), where the Lorentzian projection is exactly $C$-null and $\|v\|^{2} = 3|C(v)| = 9$.
